\section{Luồng xử lý câu hỏi bằng RAG}

Luồng xử lý bằng RAG được sử dụng khi câu hỏi cần khai thác thông tin từ tài liệu tuyển sinh hoặc khi Rasa không đủ tự tin để trả lời.

Các trường hợp thường dùng RAG:
\begin{itemize}
    \item Câu hỏi dài, phức tạp.
    \item Câu hỏi cần tổng hợp nhiều nguồn.
    \item Câu hỏi liên quan quy chế/thông báo/đề án tuyển sinh.
    \item Câu hỏi chưa có trong tập intent.
    \item Câu hỏi cần trích dẫn nguồn.
\end{itemize}

\subsection*{a. Indexing Flow}

Indexing Flow là quá trình xử lý tài liệu đầu vào để tạo kho tri thức.
\begin{enumerate}
    \item Tải tài liệu tuyển sinh lên hệ thống.
    \item Lưu tài liệu gốc vào Document Store.
    \item Trích xuất, làm sạch và chuẩn hóa văn bản.
    \item Chia tài liệu thành các đoạn nhỏ và gắn metadata.
    \item Sinh embedding cho từng đoạn.
    \item Lưu embedding vào Vector Store.
    \item Trích xuất thực thể -- quan hệ quan trọng vào Knowledge Graph (nếu cần).
\end{enumerate}

\begin{figure}[H]
    \centering
    \includegraphics[width=\textwidth,height=0.8\textheight,keepaspectratio]{\detokenize{content/source/RAG/flows/kmarag_indexing_flow.png}}
    \caption{Luồng Indexing Flow của phân hệ RAG}
    \label{fig:rag_indexing_flow_chapter3}
\end{figure}

\subsection*{b. Query Flow}

Query Flow là quá trình truy hồi và sinh câu trả lời khi người dùng đặt câu hỏi.
\begin{enumerate}
    \item Câu hỏi được chuyển từ Rasa/Action Server sang RAG.
    \item Chuẩn hóa câu hỏi và chuyển thành embedding.
    \item Vector Store tìm đoạn tài liệu liên quan.
    \item Document Store trả về nội dung đầy đủ.
    \item Có thể bổ sung ngữ cảnh bằng Knowledge Graph.
    \item Xếp hạng lại đoạn tài liệu, tạo prompt chứa câu hỏi + bằng chứng.
    \item Mô hình sinh ngôn ngữ tạo câu trả lời.
    \item Hệ thống gắn nguồn trích dẫn và trả về người dùng.
\end{enumerate}

\begin{figure}[H]
    \centering
    \includegraphics[width=\textwidth,height=0.8\textheight,keepaspectratio]{\detokenize{content/source/RAG/flows/kmarag_query_flow.png}}
    \caption{Luồng Query Flow của phân hệ RAG}
    \label{fig:rag_query_flow_chapter3}
\end{figure}
